% Mathematical proof component.  This file is intended to be included in the
% parent report after amsmath, amssymb, amsthm, longtable and booktabs are loaded.
\section{Closing the exact source-to-command interface}
\label{sec:source-command-complete}

This section establishes the intermediate interface between the literal
nine-register source and the abstract AX11H or QAX11H command scheduler.
It does not assume arithmetic meanings for named source macros.  Instead,
the supplied checker independently generates primitive-instruction templates,
matches every instruction and outgoing edge of each new source, and checks
the entire finite AX expression.  The proofs below establish the templates'
arithmetic meanings and exact clocks uniformly in $\mathrm{iPRA}$.
This is a written mathematical proof of weak-arithmetic provability, not
a newly kernel-checked proof.  The independent issue of calibration to PA
refutations is not used here and is not settled by this section.

\subsection{Primitive semantics and the complete finite binding}

A source configuration consists of a control $p$ and the nine natural
registers
\[
 (N,T,W,A,D,H,X,Y,L)=(r_0,r_1,r_2,r_3,r_4,r_5,r_6,r_7,r_8).
\]
An instruction $\mathrm{inc}(r,p')$ increments $r$ and goes to $p'$.
An instruction $\mathrm{dec}(r,p_0,p_+)$ goes to $p_0$ when $r=0$;
otherwise it decrements $r$ and goes to $p_+$.  A jump changes only the
control.  A halt instruction goes to $\operatorname{Stop}(\vec r)$,
retaining every register.  Stopped configurations are fixed points.
Use any fixed primitive recursive coding of finite tuples; these
instructions define a total primitive recursive function
$\operatorname{SrcStep}$.

The independently generated templates match all 3,671 instructions of
each source, including register arguments and both decrement destinations.
The checker also re-expands the supplied named syntax and resolves its
control tree independently by subtree intervals.  Its source-level ports
are as follows; intervals are half open.
\begin{center}
\begin{tabular}{rl}
\toprule
Port or interval&Operation\\
\midrule
$[0,10)$&Destructive comparison of $H,X$; equality goes to 10, inequality to 11\\
$10$&Halt, retaining all registers\\
$[11,25)$&Unpair $W$, then transfer its first projection into $H$\\
$[25,34)$&Fresh or positive-residual fetch prefix\\
$[34,48)$&Unpair $L$, then transfer its first projection into $A$\\
$48,49$&Two dispatch decrements of $A$\\
$[50,3572)$&The 190-token AX body, including its final extra push\\
$[3572,3639)$&Read the next field into $D$ and construct $I(T,D)$\\
$[3639,3642)$&Transfer $T$ into $X$\\
$[3642,3652)$&Compare $H,X$; equality goes to 3652, inequality to 3659\\
$[3652,3659)$&Copy $D$ to $T$ and restore $D$\\
$3659$&Jump to fetch port 25\\
$[3660,3662)$&Clear $T$\\
$[3662,3669)$&Copy $H$ to $T$ and restore $H$\\
$3669,3670$&Increment $X$, then jump to root port 0\\
\bottomrule
\end{tabular}
\end{center}
Dispatch at 48 sends zero directly to 3660.  Dispatch at 49 sends zero
to 3572; its positive branch enters AX at 50.  Thus the three cases are
the tags 0, 1, and $c+2$.

For precision, the following are the nontrivial arithmetic templates.
All indices in a template are relative to its entry; $e$ is its ordinary
exit.  The letters $D_r(a,b)$, $I_r(a)$, and $J(a)$ denote a decrement,
increment, and jump.  The comparison has two supplied exits $e_=,e_\ne$.
These lists specify the primitive programs, rather than invoking a macro
name as an arithmetic assumption.
\begin{align*}
\mathcal U_r={}&[
 D_r(8,1), I_Y(2),D_X(4,3),J(7),D_Y(7,5),I_X(6),\\
 &\hspace{12mm}J(4),J(0),D_Y(e,9),I_r(10),J(8)],\\
\mathcal P={}&[
 D_T(4,1),I_Y(2),I_X(3),J(0),D_X(e,5),I_Y(6),\\
 &\hspace{12mm}D_X(9,7),I_T(8),J(6),J(0)],\\
\mathcal C={}&[
 D_H(4,1),D_X(3,2),J(0),J(5),D_X(e_=,5),\\
 &\hspace{12mm}D_H(7,6),J(5),D_X(9,8),J(5),J(e_\ne)] .
\end{align*}
The ordinary exits of $\mathcal U_r$ and $\mathcal P$ are respectively
11 and 10.  The remaining templates are clear
$[D_r(e,1),J(0)]$, transfer
$[D_r(e,1),I_s(2),J(0)]$, duplicate
$[D_r(e,1),I_s(2),I_t(3),J(0)]$, and multiplication by four
$[D_Y(e,1),I_T(2),I_T(3),I_T(4),I_T(5),J(0)]$.
Their ordinary exits are their lengths.  Distinct register names are used
in every instantiated transfer and duplicate.

\subsection{Exact arithmetic clocks and their weak induction proofs}

Write
\[
 \pi(a,b)=\operatorname{tri}(a+b)+b,\qquad
 \operatorname{tri}(s)=\sum_{j<s}(j+1).
\]
Define $\operatorname{advance}(a,b)=(a-1,b+1)$ for $a>0$, and
$\operatorname{advance}(0,b)=(b+1,0)$.  Define
$\operatorname{decode}(w)=\operatorname{advance}^{w}(0,0)$.
All these functions use primitive recursion.  Direct arithmetic gives
$\pi(\operatorname{advance}(a,b))=\pi(a,b)+1$.  Quantifier-free
induction therefore proves $\pi(\operatorname{decode}(w))=w$.
The intervals $[\operatorname{tri}(s),\operatorname{tri}(s)+s]$ are
disjoint and successive; within one such interval $b$ is unique.
This proves injectivity of $\pi$ using quantifier-free arithmetic and
decidable comparisons, and hence
$\operatorname{decode}(\pi(a,b))=(a,b)$, including $(a,b)=(0,0)$.

Put $R(a,b)=5$ for $a>0$ and $R(0,b)=3b+8$, and define
\[
 \Lambda(w)=1+\sum_{j<w}R(\operatorname{advance}^{j}(0,0)),\qquad
 U(w)=\Lambda(w)+3\operatorname{tl}(w)+1,
\]
where $(\operatorname{hd}(w),\operatorname{tl}(w))=
\operatorname{decode}(w)$.  With $r=w$ and $X=Y=0$, the template
$\mathcal U_r$ terminates with
\[
 (r,X,Y)=(\operatorname{tl}(w),\operatorname{hd}(w),0)
\]
in exactly $U(w)$ primitive steps.  Indeed, each round at offset 0
consumes one unit of $r$.  Its two branches advance $(X,Y)$ in exactly
$R(X,Y)$ steps.  The last zero test costs one step, and the final transfer
of $Y$ into $r$ costs $3\operatorname{tl}(w)+1$.
The invariant after $j\le w$ rounds is
\[
 r=w-j,\quad(X,Y)=\operatorname{advance}^{j}(0,0),
 \quad\text{elapsed}=\sum_{i<j}R(\operatorname{advance}^{i}(0,0)).
\]
All register equalities and the corresponding iterated-step equality are
primitive recursive assertions.  The induction on $j$ is quantifier-free;
the guard $j\le w$ can itself be expressed by a Boolean characteristic
function.  This proof does not require an unbounded search for a decoder
exit.  A useful bound is $\Lambda(w)\le w(3w+8)+1$.

With $(T,X,Y)=(t,a,0)$, the pairing template $\mathcal P$ terminates
with $(T,X,Y)=(0,0,\pi(a,t))$ in
\[
 P(t,a)=4t+2+\sum_{j<a+t}(7j+5)                       \tag{P}
\]
steps.  Its first transfer costs $4t+1$ and reaches
$(T,X,Y)=(0,a+t,t)$.  Each subsequent round takes $X=s+1$ to $s$,
adding $s+1$ to $Y$, in $7s+5$ steps: subtract one and increment $Y$;
transfer the remaining $s$ units of $X$ to $T$; jump back; duplicate
those $s$ units from $T$ into $X,Y$.  The final zero test adds one.
Quantifier-free induction on the number of rounds proves the triangular
sum and the exact step identity.  All unrelated registers are parameters
and remain unchanged.

The clear, transfer, and duplicate templates have respective costs
$2r+1$, $3r+1$, and $4r+1$, with the indicated source register cleared.
Transfer adds its old value to its destination; duplicate adds it to
both destinations.  Induction on the original counter proves the exact
endpoints.  A duplicate from $r$ into $T,X$, followed by transfer from
$X$ back into $r$, copies $r$ additively into $T$ and restores $r$ in
$7r+2$ steps when $X=0$.  The multiplication template sends
$(T,Y)=(t,y)$ to $(t+4y,0)$ in $6y+1$ steps.  These proofs have linear
quantifier-free counter invariants.

The comparison template clears both $H,X$, preserves other registers,
and selects $e_=$ exactly when the original values agree.  Its exact cost is
\[
 C(h,x)=
 \begin{cases}
  3h+2,&h=x,\\
  3x+2,&h<x,\\
  2h+x+4,&h>x.
 \end{cases}                                           \tag{C}
\]
Each joint positive round subtracts one from both operands in three
steps.  The two terminal cases either exit in two steps or use the
displayed clear loops.  Equivalently,
$C(0,0)=2$, $C(0,x+1)=3x+5$, $C(h+1,0)=2h+6$, and
$C(h+1,x+1)=C(h,x)+3$.
Induction on the number of joint rounds, bounded by $\min(h,x)$,
establishes both the destination and clock.  In particular, equality
and inequality are handled for all natural operands.

Each of these proofs can be written as a quantifier-free equality for
the iterated primitive step function, with the parameters fixed at the
start of the induction.  When a remaining counter decreases, use the
number of completed rounds as the induction variable.  Bounded Boolean
conditions account for its range.  Thus no induction on an unbounded
universal-input assertion is hidden in the argument.

\subsection{Derived macros, all tied to primitive templates}

The checker expands the following compositions into the preceding literal
templates before comparing the source instruction array.  Their meanings
are therefore consequences of those templates.  Scratch registers are
$X=Y=0$ at ordinary macro boundaries.

Push is transfer $T$ to $X$, transfer $W$ to $T$, pairing $\mathcal P$,
and transfer $Y$ to $W$.  It sends $(T,W)=(t,w)$ to
$(0,\pi(t,w))$, with exact cost
\[
 B(t,w)=3t+3w+P(w,t)+3\pi(t,w)+3.
\]
CopyD, CopyH, and CopyA are Push followed by the corresponding
copy-and-restore template.  Thus they preserve the old visible top on the
stack and place the chosen register in $T$.

Cons is $\mathcal U_W$, then $\mathcal P$, then multiplication by four.
It sends $(T,W)=(t,w)$ to
\[
 (I(\operatorname{hd}(w),t),\operatorname{tl}(w)),\qquad
 I(a,b)=4\pi(a,b),
\]
in $U(w)+P(t,\operatorname{hd}(w))+6\pi(\operatorname{hd}(w),t)+1$
steps.  ReadH is $\mathcal U_L$ followed by transfer $X$ to $H$; it
adds $\operatorname{hd}(L)$ to $H$ and replaces $L$ by its tail.
ReadDCons is $\mathcal U_L$, transfer $X$ to $D$, CopyD, then Cons.
It therefore sends
\[
 (T,D,L)=(t,d,\pi(x,\ell))
 \longmapsto (I(t,d+x),d+x,\ell),
\]
preserving the original $W$ and $H$.

CopyGuard duplicates $A$ into $Y,X$, clearing $A$.
CompareGuard is: compare $H,X$; clear $T$ on equality; transfer $Y$
to $X$; transfer $D$ to $H$; compare $H,X$ again; clear $T$ on equality.
Thus on entry $(T,D,H,X,Y,A)=(g,D,H,a,a,0)$ it returns
\[
 (T,D,H,X,Y,A)=
 \begin{cases}(g,0,0,0,0,0),&a\ne H\text{ and }a\ne D,\\
 (0,0,0,0,0,0),&\text{otherwise}.
 \end{cases}
\]
The optional clear loops are charged exactly when their equality branch
is taken.  The actual primitive comparison exits are checked, including
the branch that skips each clear.  This supplies the previously assumed
guarded-cleanup meaning.

All costs compose by addition, evaluated at successive arithmetic stores.
The class \texttt{Arithmetic} in the companion script is an explicit
implementation of these endpoint and clock functions.  Its assertions
record preconditions: the mathematical PR functions are the displayed
arithmetic clauses, with arbitrary values permitted outside their claimed
domains.  The source calls below satisfy every precondition.  No duration
is obtained by waiting for an unspecified exit.

\subsection{The complete AX arithmetic expression}

Set $M(a,b)=I(a,b)+1$, $U(a,B)=I(B,a)+2$, $E(a,b)=I(a,b)+3$,
$\neg B=I(B,1)$, $B\wedge C=\neg I(B,\neg C)$, and
$\exists a\,B=\neg U(a,\neg B)$.  On AX entry put $D=d+x$ and
$H=h+y$.  Define the six heads
\begin{align*}
 h_1&=I(\exists D\,U(D,H),H),\\
 h_2&=I(U(D,U(a,H)),U(a,U(D,H))),\\
 h_3&=I(M(H,D),I(E(H,a),M(a,D))),\\
 h_4&=I(E(H,D),I(M(a,H),M(a,D))),\\
 h_5&=\exists D\,E(H,D),\\
 h_6&=I(U(D,a),I(U(D,I(a,H)),U(D,H))).
\end{align*}
For QAX11H only, replace $h_6$ by
$I(U(D,I(a,H)),I(U(D,a),U(D,H)))$.
Also put
\[
 \rho=I(I(I(H,d),D),I(I(D,H),I(a,H))),\qquad
 \Gamma=I(a,U(a,a)).
\]
The gate expression is specified without any semantic assumption about
the underlying calculus:
\begin{align*}
 p&=I(a,M(a,D)),&v&=I(M(H,D),U(a,M(H,D))),\\
 s&=U(D,I(M(H,D),p)\wedge I(M(a,D),a)),\\
 q&=U(a,I(s,\neg p)),&g&=\neg I(v,q).
\end{align*}
Write $G=g$ if $a\ne D,H$, and $G=0$ otherwise.

For every input $n,t,w,a,d,h,x,y,z,\ell$, the complete body at
port 50 starts with register file
\[
 (n,t,w,a,d,h,0,0,\pi(y,\pi(x,\pi(z,\ell))))
\]
and terminates at port 3572 with
\[
 \bigl(n,0,[G,\Gamma,h_6,h_5,h_4,h_3,h_2,h_1,
          \rho,U(d,h),t;w],0,0,0,0,0,\pi(z,\ell)\bigr).       \tag{AX}
\]
Here $[b_0,\ldots,b_k;w]=\pi(b_0,\pi(b_1,\ldots\pi(b_k,w)\ldots))$.

To verify this finite algebraic identity, apply the just-proved macro
equations in the exact 190-token order emitted by the literal grammar.
The companion's \texttt{symbolic\_AX} performs these substitutions with
independent formal variables and compares all nine output registers with
(AX).  In particular, it retains both equality tests symbolically as
zero-on-equality operations, rather than selecting a favorable guard
branch.  Its expected head expressions are the formulas above.
Every primitive instruction underlying every token has independently
matched its template, and every ordinary token has zero scratch at its
entry and exit.  The only exception is CopyGuard immediately followed by
CompareGuard, whose stated intermediate precondition is checked explicitly.
Consequently the finite comparison establishes an identity by congruence
of primitive recursive functions.  It is not a bounded numerical test of
the AX arithmetic.

Let $(F_j(\vec r),K_j(\vec r))$ be the endpoint and positive cost of
token $j$ just proved.  Define $\vec r_0=\vec r$,
$\vec r_{j+1}=F_j(\vec r_j)$, and
$K_{\mathrm{AX}}(\vec r)=\sum_{j<190}K_j(\vec r_j)$.
This fixed finite composition is an explicit PR clock for (AX).
Its correctness follows by 190 applications of the iterated-step
composition identity, itself proved by quantifier-free induction.

\subsection{The changed MP and AX fallthrough}

From port 3572 with file $(n,t,w,0,d,h,0,0,\pi(x,\ell))$,
ReadDCons constructs $q=I(t,b)$ with $b=d+x$.
The transfer at 3639 moves this $q$ from $T$ into $X$.
The comparison at 3642 clears $H,X$.  Equality enters the seven-site
copy-and-restore tail, which writes $b$ into the now-zero $T$ and preserves
$D=b$.  Inequality skips that tail.  The final jump at 3659 returns to
port 25.  The exact output is therefore
\[
 (n,\,\mathbf{if}\ q=h\ \mathbf{then}\ b\ \mathbf{else}\ 0,
     \,w,0,b,0,0,0,\ell).                                  \tag{MP}
\]
There is no stack push in the success tail.  Its exact additional cost
after ReadDCons is
\[
 (3q+1)+C(h,q)+\mathbf 1_{q=h}(7b+2)+1.
\]
The input to ReadDCons and hence its cost are fully specified, so this
gives the PR function $K_{\mathrm{MP}}$ implemented by
\texttt{Arithmetic.mp\_tail}.

Applying this result to the output (AX) consumes the additional field
$z$.  It has $t=d=h=0$, so $b=z$ and $q=I(0,z)$.
The elementary identity $I(0,z)=0\leftrightarrow z=0$ follows directly
from Cantor pairing.  If $z=0$, the successful copy writes zero; if $z>0$,
comparison fails and $T$ is also zero.  Thus the complete AX command
returns $(T,W,D,H)=(0,[G,\Gamma,h_6,\ldots,\rho,U(d,h),t;w],z,0)$
with residual $\ell$, in exactly
$K_{\mathrm{AX}}+K_{\mathrm{MP}}$ steps, the latter evaluated at the
actual AX-body output.  This handles both zero and nonzero $z$.

\subsection{The cached-only Shift and exact halted file}

From port 3660, clear $T$, copy and restore $H$, increment $X$, and
jump to root port 0.  The file now has $(T,H,X)=(h,h,1)$ and unchanged
$N,W,D,L$; $A=Y=0$.  This preparation costs
\[
 (2t+1)+(7h+2)+2.
\]
The root comparison costs $C(h,1)$.  If $h=1$, it reaches port 10
with $H=X=0$ and $T=1$, and the halt instruction costs one more step.
The exact terminal configuration is
\[
 \operatorname{Stop}(n,1,w,0,d,0,0,0,\ell).                 \tag{SH}
\]
In particular, the preserved stack and registers are not silently erased.
If $h\ne1$, the comparison goes to 11, unpairs $W$, and transfers its
first projection into the now-zero $H$.  The exact live output at 25 is
\[
 (n,h,\operatorname{tl}(w),0,d,\operatorname{hd}(w),0,0,\ell).
\]
The extra cost in this branch is $U(w)+3\operatorname{hd}(w)+1$.
These statements include $w=0$, $h=0$, arbitrary old $t$, and the case
$t=h=1$.  The Shift halt test depends only on the old cache $h$.

\subsection{Fetch, total parsing, and the composed clock}

At port 25, let the live file be $(n,t,w,0,d,h,0,0,\ell)$.
The nine-site fetch prefix has the exact effect
\[
 (n,\ell)\longmapsto(n',s),\qquad
 (n',s)=
 \begin{cases}(n+1,n),&\ell=0,\\(n,\ell-1),&\ell>0.
 \end{cases}                                             \tag{F}
\]
In the fresh case the source first duplicates $n$ into $L,X$, clearing
$N$, then increments $N$ while draining $X$.  Its cost is
$1+(4n+1)+(3n+2)=7n+4$.  In the positive-residual case it decrements
$L$ and jumps, costing two steps.  Thus the fresh counter increments
\emph{when the candidate is selected}, not when its continuation finishes.

The decoder at 34 and transfer to $A$ cost
$U(s)+3\operatorname{hd}(s)+1$.  Write
$(a,v)=\operatorname{decode}(s)$.  The two dispatch instructions
give the following total parser:
\[
 \operatorname{parse}(s)=
 \begin{cases}
 (S,v),&a=0,\\
 (\mathrm{MP}_x,v_1),&a=1,\\
 (\mathrm{AX}_{a-2,x,y,z},v_3),&a\ge2.
 \end{cases}                                             \tag{Parse}
\]
For MP use $(x,v_1)=\operatorname{decode}(v)$. For AX use the three
successive field decodings
\[
 (y,v_1)=\operatorname{decode}(v),\qquad
 (x,v_2)=\operatorname{decode}(v_1),\qquad
 (z,v_3)=\operatorname{decode}(v_2).
\]
Dispatch costs one step for $S$ and two otherwise.  The MP and AX field
decodings are charged in their respective command clocks, not twice.
Every natural seed has exactly one of these interpretations.  No syntactic
well-formedness assumption is imposed on a candidate.  Reconstructing each
decoded seed by Cantor pairing supplies the input forms in (AX) and (MP).
The identities hold when any or every field is zero.

Define the exact one-step clock $\theta(n,\ell,t,w,d,h)$ by adding:
the appropriate fetch-prefix cost in (F); the decoder/selector cost;
the appropriate one- or two-step dispatch cost; and the command clock
given above, evaluated at that actual selected register file.
This is the explicit algorithm \texttt{command\_clock}.  Its only loops
are bounded sums or iterations defining Cantor decoding and costs, and the
fixed 190-token AX composition.  It is therefore primitive recursive.
The preceding template equations prove that all scratch preconditions
hold for every live six-tuple of natural numbers.  Every component costs
a nonnegative number of steps, and fetch alone costs at least two, so
$\theta\ge1$ uniformly.  The clock includes the Shift return through
the root guard and every command's return to fetch.

\subsection{The completed iPRA simulation and its consequence}

Define the scheduler $R$ on live states $(n,\ell,t,w,d,h)$ by (F),
(Parse), and the stated arithmetic command effects.  A halting Shift
goes to the data-carrying terminal state
$\operatorname{Done}(n',\ell',w,d)$, where $\ell'$ is its parsed
residual.  Done states are fixed points.  This enriches the representation
of an abstract halt without changing which computations halt.
Define
\begin{align*}
 J(n,\ell,t,w,d,h)&=(25;(n,t,w,0,d,h,0,0,\ell)),\\
 J(\operatorname{Done}(n,\ell,w,d))&=
       \operatorname{Stop}(n,1,w,0,d,0,0,0,\ell).
\end{align*}
Set $\theta=1$ on Done states.  The halted source is absorbing and
retains precisely the file in (SH).

\begin{theorem}[Exact source-to-scheduler binding]
For either of the two supplied literal source instruction arrays,
$\mathrm{iPRA}$ proves, uniformly in every scheduler state $c$,
\[
 \theta(c)\ge1,\qquad
 \operatorname{SrcStep}^{\theta(c)}(J(c))=J(R(c)),\qquad
 h_{\mathrm{Src}}(J(c))=h_R(c).
\]
In particular, starting at $c_0=(2,0,0,0,0,0)$, it proves
\[
 \operatorname{Halt}(\operatorname{Src},J(c_0))
 \ \longleftrightarrow\ \operatorname{Halt}(R,c_0).
\]
\end{theorem}
\begin{proof}
For live states, compose the exact fetch, dispatch, and command equations
proved above.  Their cases cover all natural inputs; all intervening ports
and complete register files agree.  For Done states the equation follows
from source absorption and the chosen one-step clock.  The displayed
halting flags agree by the two clauses defining $J$.

For completeness, put $B(0)=0$ and
$B(k+1)=B(k)+\theta(R^k(c_0))$.  Quantifier-free induction yields
$B(k)\ge k$ and
$\operatorname{SrcStep}^{B(k)}(J(c_0))=J(R^k(c_0))$.
A scheduler halt therefore gives a source halt at the displayed PR time.
Conversely, a source halt at time $u$ remains a halt at time $B(u)\ge u$,
by source absorption.  The endpoint equality and halt-flag agreement show
that $R^u(c_0)$ is halted.  Both implications are constructive and use no
unbounded search or classical principle.
\end{proof}

This theorem supplies the previously separate primitive-source/command
interface.  Combining it with a completed iPRA physical/source simulation
and the abstract command-preservation theorem identifies literal-machine
halting with numerical-calculus refutation.  It does not identify that
calculus with PA without the separate logical proof translations.

\subsection{Reproduction and scope of the finite receipts}

The checker uses Python's standard library only.  Run it separately on
\texttt{machines/PA194AX11H} and \texttt{machines/PA194QAX11H}, specifying
an output directory.  Each receipt records the four exact input hashes,
the complete instruction/edge match, independent tree resolution,
190 macro tokens, and the free-algebra identity.  It also records 882
representative exact local-duration checks and 324 composed Shift/MP
checks per family.  Those numerical checks detect implementation mistakes
in clocks; the uniform claims follow from the counter-induction proofs
and finite source binding above.  No practical execution bound or
sampling cutoff is asserted for the enormous full AX arithmetic.
